% Copyright (c) 2026 Geovana Neves. Licensed under CC BY 4.0 (https://creativecommons.org/licenses/by/4.0/). See LICENSE-AND-AUTHORSHIP.md.
%
% workspaces/text/01-tree-and-inputs.tex
%
% Parts 1 and 2 of guide 2: the tree and the inputs.

\section{The tree}

\begin{frame}[fragile]{One folder holds the campaign}
\begin{columns}[T]
\begin{column}{0.56\textwidth}
\begin{lstlisting}
campaign/
  init.toml              the solver of this workspace
  inputs/                yours: everything a run reads
    matrices/<name>.fs
    setups/s<id>.toml
    references/r<id>.toml
    pproc/p<id>.toml
    mesh/
  sims/                  one folder per simulation
    sim_<POL>/datapoints/DP-<point name>/
  post/<matrix>/         plan receipt, tables
  archive/<stamp>/       earlier evidence
  runs.json              the record of every point
\end{lstlisting}
\end{column}
\begin{column}{0.4\textwidth}
\begin{itemize}\small
  \item The tree is the same for every solver.
  \item A study moves to another machine by copying or syncing the folder.
  \item Nothing lives in a database or in a hidden folder of your account.
\end{itemize}
\begin{takeaway}If it is not in the folder, it is not part of the
campaign.\end{takeaway}
\end{column}
\end{columns}
\end{frame}

\begin{frame}{Who writes what}
\relax{\cfdbbodysize
\begin{tabularx}{\linewidth}{@{}P{0.27\linewidth}P{0.27\linewidth}R@{}}
\toprule
\textbf{Where} & \textbf{Written by} & \textbf{Yours to edit?} \\
\midrule
\code{init.toml} & \code{cfd-workspace init} & no \\
\code{inputs/} & you & yes, and only here \\
\code{sims/} & \code{cfd-matrix run} & no \\
\code{post/<matrix>/} & \code{plan}, \code{collect}, \code{post} & no:
  rebuilt each time \\
\code{runs.json} & \code{run}, \code{collect} & no \\
\code{marks.json} & \code{mark-failed}, \code{mark-converged} & no \\
\code{archive/<stamp>/} & any command that replaces something & no \\
\code{logs/} & \code{run}, \code{collect}, \code{post} & no \\
\bottomrule
\end{tabularx}}
\vspace{2mm}
\begin{takeaway}The files under \code{post/} are derived: deleting them loses
nothing, since \code{collect} rebuilds them from the records.\end{takeaway}
\end{frame}

\begin{frame}[fragile]{Creating a workspace}
\begin{lstlisting}
cfd-workspace init campaign --solver openfoam
\end{lstlisting}
\begin{itemize}
  \item Creates the folder if needed, \code{init.toml}, the input folders,
        \code{sims/}, \code{post/} and \code{archive/}.
  \item \code{--solver} is required. The names are lower case: \code{su2},
        \code{openfoam}, \code{pyfr}, \code{xfoil}, \code{qprop},
        \code{mock}, \code{flightstream}.
  \item A folder that already is a workspace is refused: a workspace keeps
        its solver.
  \item It also writes \code{inputs/.gitignore}, so the file that holds this
        machine's solver paths stays out of a repository.
\end{itemize}
\begin{cfdbpitfall}[One solver per workspace]
To compare two solvers on the same matrix, make two workspaces and copy the
matrix. The columns mean the same in both.
\end{cfdbpitfall}
\end{frame}

\section{The inputs}

\begin{frame}{The matrix columns}
\relax{\cfdbtablesize
\begin{tabularx}{\linewidth}{@{}P{0.25\linewidth}R@{}}
\toprule
\textbf{Column} & \textbf{Meaning} \\
\midrule
\code{POL} (required) & the simulation; its folder is \code{sims/sim_<POL>/}.
  Unique per matrix \\
\code{RUN} (required) & \code{1} plans and runs the row, \code{0} keeps it
  without running it \\
\code{SET}, \code{REF}, \code{PPROC} (required) & the setup, reference and
  post file of the row \\
\code{FLIGHT_CONDITION}, \code{SWEEP_VALUES} & the flow state and attitude;
  the values of the one key written \code{sweep} \\
\code{GEOMETRY} & the mesh or geometry file in \code{inputs/mesh/} \\
\code{AIRCRAFT}, \code{CONFIGURATION}, \code{DESCRIPTION} & your labels; they
  reach the plan listing and the summary \\
\code{NCPUS}, \code{WALLTIME} & what a cluster job asks for; \code{WALLTIME}
  carries its unit: \code{90s}, \code{240m}, \code{4h}, \code{1d} \\
\code{SOLVER_BUILD} & the solver version the row runs \\
\code{VAR_NAMES_VALUES} & per-row overrides of setup keys \\
\code{HIDDEN}, \code{SYMMETRY}, \code{SYMMETRY_LOADS}, \code{WORKFLOW} &
  kept so one layout serves every solver; the library solvers carry them or
  accept \code{-} only \\
\bottomrule
\end{tabularx}}
\vspace{1mm}
{\small An optional column may be left out of the header. A column the layout
does not know is refused before anything is planned.\par}
\end{frame}

\begin{frame}{The flight condition: the air}
\begin{columns}[T]
\begin{column}{0.5\textwidth}
\relax{\cfdbtablesize
\begin{tabularx}{\linewidth}{@{}P{0.24\linewidth}P{0.2\linewidth}R@{}}
\toprule
\textbf{Key} & \textbf{Unit} & \textbf{Role} \\
\midrule
\code{MACH} & - & velocity \\
\code{TASmps} & m/s & velocity \\
\code{REmi} & millions & density, from the Reynolds number \\
\code{ALTFT} & ft & pressure and temperature \\
\code{dISA} & K & temperature offset \\
\code{RHOkgm3} & kg/m3 & pins the density \\
\code{MUPas} & Pa s & pins the viscosity \\
\code{ASMPS} & m/s & pins the speed of sound \\
\code{TK} & K & pins the temperature \\
\code{PPA} & Pa & pins the pressure \\
\bottomrule
\end{tabularx}}
\end{column}
\begin{column}{0.46\textwidth}
\begin{itemize}\small
  \item A key carries its unit in its name.
  \item State exactly one of \code{MACH} or \code{TASmps}.
  \item No altitude means sea level of the standard atmosphere [5]; no
        \code{dISA} means zero.
  \item The viscosity follows Sutherland's law [6] unless pinned.
  \item \code{REmi} is measured against the \code{length} of the reference
        file.
  \item \code{REmi} beside \code{RHOkgm3} is refused: both fix the density.
\end{itemize}
\end{column}
\end{columns}
\end{frame}

\begin{frame}[fragile]{The flight condition: attitude and sweep}
\begin{lstlisting}
FLIGHT_CONDITION                              | SWEEP_VALUES
MACH:0.85, REmi:5, ALPHA:sweep, BETA:0        | 0,1,2,2.5,3
TASmps:40, ALTFT:5000, ALPHA:4                | -
\end{lstlisting}
\begin{itemize}\small
  \item \code{ALPHA} and \code{BETA} are in degrees. Keys match without
        regard to case; a key written twice is refused.
  \item Exactly one key may say \code{sweep}. Each value of
        \code{SWEEP_VALUES} is one \textbf{point} of the simulation.
  \item The point's folder name is built from its flight condition, so the
        same operating point has the same name in every workspace.
  \item A solver may read keys of its own, such as the rotation rate of a
        propeller. Guide 3 lists them.
\end{itemize}
\begin{cfdbpitfall}[A bare Mach number]
For some solvers a row that states only \code{MACH} and angles keeps the flow
state of its setup file. State the air in the row when you want the row to
decide it.
\end{cfdbpitfall}
\end{frame}

\begin{frame}[fragile]{Overriding a setup key in one row}
\begin{lstlisting}
VAR_NAMES_VALUES
iterations: 4000 / solver.threads: 4
\end{lstlisting}
\begin{itemize}
  \item \code{KEY: value} pairs separated by \code{ / }.
  \item Each key replaces the key of the same name in the row's setup file.
        A dotted key reaches inside a table: \code{solver.threads} is
        \code{threads} in \code{[solver]}.
  \item Use the name the setup file uses.
  \item A key that is a matrix column, such as \code{NCPUS}, is refused
        here: state it in its column.
\end{itemize}
\begin{takeaway}One setup file and a column of overrides is easier to review
than ten setup files that differ in one line.\end{takeaway}
\end{frame}

\begin{frame}{The mesh is stored once}
\begin{itemize}
  \item The \code{GEOMETRY} cell names a file in \code{inputs/mesh/}: a file
        name, or a stem that names exactly one file there.
  \item A mesh is \textbf{never copied} into a point folder. Each point
        folder holds a link to the file in \code{inputs/mesh/}.
  \item The link is relative, so it still resolves after the whole workspace
        moves to another machine.
  \item The mesh's hash is recorded with the plan and checked before the
        run.
\end{itemize}
\begin{cfdbpitfall}[Do not edit a mesh in place]
Points already run are linked to that file. A changed mesh is a new file with
a new name.
\end{cfdbpitfall}
\end{frame}

\begin{frame}[fragile]{Where the solver is: executables}
\begin{columns}[T]
\begin{column}{0.5\textwidth}
\begin{lstlisting}[style=cfdbtoml]
# inputs/executables.toml
[su2]
"8.3.0" = "/opt/su2-8.3.0/bin/SU2_CFD"
\end{lstlisting}
\begin{lstlisting}[style=cfdbtoml]
# inputs/executables.local.toml
[su2]
"8.3.0" = "D:/solvers/su2-8.3.0/SU2_CFD.exe"
\end{lstlisting}
\end{column}
\begin{column}{0.46\textwidth}
\begin{itemize}\small
  \item The first file declares the builds a row may name in
        \code{SOLVER_BUILD}. It travels with the workspace.
  \item The second gives \textbf{this machine's} paths for those builds. It
        stays out of version control.
  \item A build no file declares refuses the run before anything is
        written.
  \item Without these files, the setup's \code{[solver]} table says where
        the solver is.
\end{itemize}
\end{column}
\end{columns}
\end{frame}
